% Generated draft: requires mathematical and source review.
\documentclass[11pt]{article}
\usepackage[margin=25mm]{geometry}
\usepackage{fontspec}
\setmainfont{Latin Modern Roman}
\usepackage{amsmath,amssymb,mathtools,unicode-math}
\setmathfont{Latin Modern Math}
\usepackage{enumitem,needspace,xurl,hyperref}
\setlength{\parindent}{0pt}
\setlength{\parskip}{5pt}
\setlength{\emergencystretch}{3em}
\newcommand{\N}{\mathbb{N}}
\newcommand{\Z}{\mathbb{Z}}
\newcommand{\Q}{\mathbb{Q}}
\newcommand{\R}{\mathbb{R}}
\newcommand{\C}{\mathbb{C}}
\newcommand{\F}{\mathbb{F}}
\newcommand{\A}{\mathbb{A}}
\newcommand{\PP}{\mathbb{P}}
\newcommand{\E}{\mathbb{E}}
\newcommand{\eps}{\varepsilon}
\newcommand{\dd}{\,\mathrm{d}}
\begin{document}
\section*{UnsolvedMath ID: 30003890}
\textbf{Source problem code:} catalogue entry
\subsection{Definitions and mathematical statement}
Let \(n\) be an integer with \(n\geq 3\), and put \(J=\operatorname{diag}(1,\ldots,1,-1)\). Define
\[
G_n=\operatorname{SU}(n,1)=\{A\in\operatorname{SL}_{n+1}(\C):A^*JA=J\},
\]
where \(A^*\) denotes the conjugate transpose. Complex hyperbolic space \(\mathbb H_{\C}^n\) is the space of negative complex lines for this Hermitian form, equipped with its complex hyperbolic metric. Its holomorphic isometry group is
\[
\operatorname{PU}(n,1)\cong G_n/Z(G_n),\qquad
Z(G_n)=\{\zeta I:\zeta^{n+1}=1\}.
\]
The full isometry group also contains antiholomorphic isometries. Thus the source's displayed identification of \(\operatorname{Isom}(\mathbb H_{\C}^n)\) with \(\operatorname{SU}(n,1)\) requires correction: the holomorphic group is a finite central quotient of \(G_n\), and the full isometry group is larger.

A subgroup \(\Gamma\leq G_n\) is a \emph{lattice} if it is discrete and \(\Gamma\backslash G_n\) has finite Haar volume. No compactness of this quotient is required. Two subgroups of the same group are \emph{commensurable} if their intersection has finite index in each.

Use the following conventional arithmeticity definition. A lattice \(\Gamma\leq G_n\) is \emph{arithmetic} if there exist a connected semisimple algebraic group \(\mathbf H\) over \(\Q\), a faithful rational matrix representation \(\mathbf H\hookrightarrow\operatorname{GL}_m\), a subgroup \(\Lambda\leq\mathbf H(\Q)\) commensurable with \(\mathbf H(\Q)\cap\operatorname{GL}_m(\Z)\), and a surjective Lie group homomorphism
\[
\phi:\mathbf H(\R)^\circ\longrightarrow G_n
\]
with compact kernel, such that \(\Gamma\) is commensurable with \(\phi(\Lambda\cap\mathbf H(\R)^\circ)\). Here the superscript \(\circ\) denotes the identity component. A lattice is \emph{nonarithmetic} if it is not arithmetic. Passing between \(G_n\) and its finite central quotient preserves the lattice existence question and arithmeticity.

\textbf{Question, interpreted dimension by dimension.} For each integer \(n\geq3\), determine whether there exists a nonarithmetic lattice in \(G_n\); equivalently, determine the set
\[
\bigl\{n\in\Z:n\geq3\text{ and there exists a nonarithmetic lattice }\Gamma\leq G_n\bigr\}.
\]
The source does not explicitly distinguish a question for every dimension from a question asking for an example in some dimension \(n\geq3\). The dimension-by-dimension interpretation above is supported by its supplied literature assessment and remains a qualification requiring review.

The supplied curation labels the problem \emph{partially solved}: it reports examples in \(\operatorname{SU}(2,1)\) and \(\operatorname{SU}(3,1)\), and reports existence in dimensions \(n\geq4\) as open. Thus its stated remaining target is the same existence question for each integer \(n\geq4\). These are supplied dated assessments, not independently verified mathematical or current-status claims.
\subsection{Short English statement}
For each complex dimension \(n\geq3\), ask whether complex hyperbolic space admits a nonarithmetic lattice, meaning a discrete group of finite covolume that does not arise from the arithmetic construction specified above. The supplied assessment reports dimension \(3\) as resolved and dimensions \(n\geq4\) as open. The source's isometry-group notation and intended quantifier require the qualifications recorded here.
\subsection{Sources}
\begin{itemize}
\item Local UnsolvedMath registry: supplied record with permanent ID \(30003890\), exact source problem code \texttt{catalogue entry}, title \emph{Nonarithmetic Lattices in Higher-Dimensional Complex Hyperbolic Space}, in the collection \emph{Oberwolfach Reports Open Problems}. The project identifies \texttt{data/unsolvedmath/problems.json} and \texttt{catalog/problems.json} as its registry files; their contents were not independently inspected for this draft.
\item Supplied primary report citation: \emph{Topologie (2019)}, Oberwolfach Reports, pp.\ 1857--1910, DOI \texttt{10.4171/owr/2018/31}, \url{https://doi.org/10.4171/owr/2018/31}. The record gives proposed year \(2018\) and does not name a proposer. Its background describes Oberwolfach Reports as documenting mathematical workshops and their open-problem discussions.
\item Supplied literature source: Parker, non-arithmetic lattices, \url{https://www.maths.ox.ac.uk/node/13849}.
\item Supplied literature source: Mostow, non-arithmetic lattice in \(\operatorname{SU}(2,1)\), \url{https://pmc.ncbi.nlm.nih.gov/articles/PMC392706/}.
\item Supplied literature source: Oberwolfach report, \url{https://doi.org/10.4171/owr/2018/31}.
\item Literature assessment date, copied exactly: \texttt{2026-08-21}. This is supplied dated curation, not independently verified. It reports nonarithmetic lattices in \(\operatorname{SU}(2,1)\) and \(\operatorname{SU}(3,1)\), and reports existence for \(n\geq4\) as open. The supplied background calls this a dated literature triage derived from OWR curation and requests independent verification of time-sensitive status and cited claims.
\item The registry supplies status \texttt{partially\_solved} and statement status \texttt{corrected\_verified}, with the explanation that mathematical content was preserved while notation, typography, and omitted definitions were repaired. These labels record the supplied editorial assessment and do not constitute independent verification here.
\end{itemize}
\end{document}
